% ── Chapter 3: Proposed Reference Framework and Analytical Evaluation ────────
% Target length: 15–18 pages | Carries construct C3
% This is the chapter the jury weighs most.
\chapter{Proposed Reference Framework and Analytical Evaluation}

% ─────────────────────────────────────────────────────────────────────────────
\section*{Introduction}
\addcontentsline{toc}{section}{Introduction}

\todo{Write chapter introduction. Signal what the jury will find: (1)~a formal artifact
with named layers, component definitions, and an integration contract — not a block
diagram with no specification; (2)~a methodical, criteria-based evaluation — not a bare
opinion table. Both attack surfaces are closed in this chapter.}

% ─────────────────────────────────────────────────────────────────────────────
\section{Requirements and Design Rationale}
\label{sec:ch3-requirements}

\todo{Derive requirements directly and traceably from RQ1--RQ3 and the ten
dual-scalability dimensions from \S\ref{sec:ch1-scalability}. Present as a requirements
table: each requirement → the RQ it addresses → the scalability dimension it serves.
This traceability is what makes the framework look engineered rather than invented.}

% Visual: requirements × RQ × dimension traceability table.

% ─────────────────────────────────────────────────────────────────────────────
\section{The Reference Framework}
\label{sec:ch3-framework}

% *** construct C3 — do not stop at a picture ***
% Specify the framework with: layered model, component definitions, integration
% contract, and numbered design invariants.

\subsection{Layered Reference Model}
\label{subsec:ch3-layers}

\todo{Define five named layers with their responsibilities:
\begin{enumerate}
    \item \textit{Ingestion \& staging layer} — heterogeneous sources, CDC, lake landing.
    \item \textit{Agile modeling layer} — DV2.0 Raw Vault + Business Vault on an open
    table format (Delta/Iceberg).
    \item \textit{Consumption/serving layer} — dimensional marts and serving views for BI.
    \item \textit{Native analytics layer} — feature pipelines, model training/serving,
    feedback loops; positioned to read from the modeling layer, not bolted past the marts.
    \item \textit{Governance \& metadata plane} — cross-cutting: lineage, schema registry,
    data catalog.
\end{enumerate}}

% Visual: layered reference architecture diagram with the governance plane cross-cutting.

\subsection{Component Definitions}
\label{subsec:ch3-components}

\todo{For each component in the five layers, state: (1)~inputs, (2)~outputs,
(3)~responsibility, (4)~which scalability dimension it serves. This is what makes
the framework a formal artifact, not a sketch.}

\subsection{The Integration Contract}
\label{subsec:ch3-contract}

\todo{This is the centerpiece — it resolves or explicitly scopes the composability tension
from \S\ref{sec:ch2-tension}. Define precisely: (1)~what the agile modeling layer
guarantees to the native analytics layer (e.g., historized, business-keyed,
lineage-stamped, feature-ready entities); (2)~how features are derived from
Satellites/Business Vault; (3)~how schema evolution in the modeling layer is propagated
without breaking trained models; (4)~how the open table format mediates both
relational-lineage discipline and ML file access. State assumptions explicitly where
you scope rather than fully solve — honesty here is strength, not weakness.}

\subsection{Design Rules and Invariants}
\label{subsec:ch3-invariants}

\todo{Enumerate numbered design rules. Numbered invariants are inexpensive to write and
read as genuine formalism — they are what a jury can probe one by one. Examples to
develop and refine:
\begin{enumerate}
    \item[R1.] Every analytical feature must trace to a historized Satellite attribute.
    \item[R2.] The native analytics layer reads exclusively from the agile modeling layer
    or the consumption/serving layer; it must never read raw staging data directly.
    \item[R3.] Schema changes in the Raw Vault must be additive (new Satellites or
    Satellite columns); destructive changes are prohibited.
    \item[R4.] Business keys must be stable and source-system-independent (hash keys).
    \item[R5.] Every data entity entering the analytics layer must carry a lineage stamp
    traceable to its source Satellite and load timestamp.
\end{enumerate}
Extend and adjust these invariants during writing.}

% ─────────────────────────────────────────────────────────────────────────────
\section{Evaluation Methodology}
\label{sec:ch3-methodology}

\todo{Before any comparison, define \emph{how} you evaluate. Short but mandatory.
Cover: (1)~method type — conceptual, criteria-based comparative assessment grounded in
literature evidence; explicitly non-empirical; (2)~the six criteria from
\S\ref{sec:ch2-criteria} restated with their rubric definitions; (3)~scoring scheme —
qualitative levels (Low/Medium/High) with the meaning of each level defined
\emph{per criterion}; (4)~evidence rule — every rating must be justified by a cited
source or a stated argument, no bare scores; (5)~threats to validity stated up front
(subjectivity, literature gaps, vendor-source bias).}

% ─────────────────────────────────────────────────────────────────────────────
\section{Comparative Study}
\label{sec:ch3-comparison}

\todo{Apply the six-criteria methodology systematically across four approaches:
\textbf{Kimball}, \textbf{Inmon}, \textbf{DV2.0}, and the \textbf{proposed framework}.
Present as an evaluation matrix (approaches × criteria). Follow with a paragraph of
justification per cell-cluster (not per cell) backed by citations.
Important: the proposed framework must \emph{not} trivially win every criterion — show
where it inherits DV2.0's implementation complexity cost. A framework that honestly
concedes weaknesses is more credible than one that claims perfection.
Anchors: BARC/Eckerson adoption figures (AnalyticsCreator 2025); Kimball \& Ross 2013;
Linstedt \& Olschimke 2015; Armbrust et al.\ 2021.}

% Visual: evaluation matrix (approaches × criteria) with rubric legend.

% ─────────────────────────────────────────────────────────────────────────────
\section{Discussion}
\label{sec:ch3-discussion}

\todo{Cover: (1)~trade-offs the framework makes and why; (2)~conditions of applicability
— when this architecture is worth its complexity (high source heterogeneity + ambition
toward predictive maturity) vs when it is overkill (stable, descriptive-only contexts);
(3)~threats to validity revisited; (4)~limitations — no empirical validation, reliance
on industry evidence for composability claims.}

% ─────────────────────────────────────────────────────────────────────────────
\section*{Conclusion}
\addcontentsline{toc}{section}{Conclusion}

\todo{Summarize C3: the layered reference framework, the integration contract, the design
invariants, and the evaluation results. Confirm which RQs are addressed and how. Bridge
to the General Conclusion, which answers all three RQs and states limitations and future
work.}
